\documentclass[]{beamer}

\usetheme{metropolis}
\metroset{progressbar=frametitle}
\usefonttheme[onlymath]{serif}

\usepackage[none]{hyphenat}
\usepackage{hyperref}
\usepackage{amsmath}
\usepackage{amsfonts}
\usepackage{amssymb}
\usepackage[normalem]{ulem}
\usepackage{graphicx}
\usepackage{tikz}
\usepackage{minted}
\author{Maxim Borisyak, Claude Code, Andrey Ustyuzhanin}
\institute{Constructor University Bremen}
\usepackage{booktabs}




\begin{document}



\title{Reviewing Scientific Articles}

\begin{frame}[plain]
\titlepage
\end{frame}

\section{The Peer Review Process}

\begin{frame}[fragile]
\frametitle{What is Peer Review?}
\textbf{Peer review}: evaluation of scientific work by experts in the same field before publication.

Occurs at:
\begin{itemize}
  \item \textbf{conference submissions} (1-3 weeks per review);
  \item \textbf{journal submissions} (months to years);
  \item \textbf{grant applications} (funding agencies);
  \item \textbf{community code review} (open source).
\end{itemize}

Reviewers are typically unpaid volunteers.

\end{frame}

\begin{frame}[fragile]
\frametitle{Goals of Peer Review}
Reviewers assess:
\begin{enumerate}
  \item \textbf{correctness}: are the claims true and properly justified?
  \item \textbf{novelty}: does this advance the field?
  \item \textbf{significance}: does this matter?
  \item \textbf{clarity}: is the work well-presented?
  \item \textbf{rigor}: is the methodology sound?
\end{enumerate}

Peer review is not:
\begin{itemize}
  \item \textbf{exhaustive verification} — reviewers check plausibility, not every calculation;
  \item \textbf{guaranteed to catch fraud} — but it catches many honest mistakes;
  \item \textbf{democratizing} — not all reviewers are equally competent.
\end{itemize}

\end{frame}

\begin{frame}[fragile]
\frametitle{Peer Review is Imperfect}
\textbf{Known limitations}:
\begin{itemize}
  \item major errors slip through (see Hwang Woo-suk, Stapel, Schön);
  \item bias exists — reviewers favor citations of their own work;
  \item reviewer selection affects outcome (expertise, competition);
  \item conflicts of interest are underreported;
  \item publication bias — negative results rarely appear.
\end{itemize}

Despite flaws, peer review remains the best mechanism we have.

\end{frame}

\begin{frame}[fragile]
\frametitle{Your Role as a Reviewer}
You are a \textbf{critical reader}, not an auditor:
\begin{itemize}
  \item catch obvious errors and inconsistencies;
  \item identify questionable claims that lack support;
  \item spot methodological red flags;
  \item assess clarity and organization;
  \item provide constructive feedback.
\end{itemize}

You are \textbf{not} expected to:
\begin{itemize}
  \item verify every calculation by hand;
  \item reproduce all experiments (though you should check design);
  \item detect all fraud (that's nearly impossible);
  \item gatekeep for competitive reasons.
\end{itemize}

\end{frame}

\section{Reading a Paper: Framework}

\begin{frame}[fragile]
\frametitle{The Reading Strategy: First Pass}
\textbf{First pass} (10-20 minutes):
\begin{itemize}
  \item title, abstract, headings, conclusion;
  \item figures and tables (captions only);
  \item determine: is this relevant to me? what's the scope?
\end{itemize}

\end{frame}

\begin{frame}[fragile]
\frametitle{The Reading Strategy: Detailed Passes}
\textbf{Second pass} (main evaluation):
\begin{itemize}
  \item introduction: problem and novelty;
  \item method: does approach make sense?
  \item results: are they convincing?
\end{itemize}

\textbf{Third pass} (if needed):
\begin{itemize}
  \item proofs or calculations;
  \item experimental design details.
\end{itemize}

\end{frame}

\begin{frame}[fragile]
\frametitle{Introduction and Motivation}
\textbf{What to check}:
\begin{itemize}
  \item \textbf{clarity of problem}: can you state what's being solved?
  \item \textbf{motivation}: why does this matter?
  \item \textbf{related work}: relevant prior art acknowledged?
  \item \textbf{claimed novelty}: what is genuinely new?
\end{itemize}

\textbf{Red flags}:
\begin{itemize}
  \item problem poorly motivated;
  \item prior work omitted or misrepresented;
  \item novelty claims without citing prior work.
\end{itemize}

\end{frame}

\begin{frame}[fragile]
\frametitle{The Method/Theory Section}
For \textbf{empirical papers}:
\begin{itemize}
  \item algorithm description clear and complete?
  \item hyperparameters fully specified?
  \item reproducible from description?
\end{itemize}

For \textbf{theoretical papers}:
\begin{itemize}
  \item definitions precise and match standard notation?
  \item assumptions clearly stated and reasonable?
  \item proofs logical and complete?
\end{itemize}

\textbf{Red flags}:
\begin{itemize}
  \item vague algorithm description;
  \item missing hyperparameters;
  \item theorems without proofs (or deferred to missing appendix).
\end{itemize}

\end{frame}

\begin{frame}[fragile]
\frametitle{Experimental Setup: Baselines and Data}
For ML/systems papers:

\textbf{Baselines}:
\begin{itemize}
  \item fair comparison? (same tuning effort)
  \item recent and from reputable implementations?
\end{itemize}

\textbf{Datasets}:
\begin{itemize}
  \item public/standard or custom?
  \item data leakage between train/test?
\end{itemize}

\end{frame}

\begin{frame}[fragile]
\frametitle{Experimental Setup: Metrics and Design}
\textbf{Metrics}:
\begin{itemize}
  \item appropriate for problem?
  \item multiple metrics or cherry-picked?
\end{itemize}

\textbf{Design}:
\begin{itemize}
  \item multiple runs with error bars?
  \item significance testing?
  \item complete ablations?
\end{itemize}

\end{frame}

\begin{frame}[fragile]
\frametitle{Results and Claims}
\textbf{What to check}:
\begin{itemize}
  \item improvements meaningful or within noise?
  \item p-values, confidence intervals reported?
  \item all runs reported or cherry-picked?
  \item conclusions match the data?
\end{itemize}

\textbf{Red flags}:
\begin{itemize}
  \item tiny improvements (0.5%) without error bars;
  \item inconsistent across datasets;
  \item "state-of-the-art" vs. outdated baselines;
  \item missing or incomplete ablations.
\end{itemize}

\end{frame}

\begin{frame}[fragile]
\frametitle{Figures and Tables: What to Look For}
\textbf{Good figures}:
\begin{itemize}
  \item clear labels, legend, caption;
  \item error bars visible;
  \item caption explains content.
\end{itemize}

\end{frame}

\begin{frame}[fragile]
\frametitle{Figures and Tables: Red Flags}
\textbf{Suspicious patterns}:
\begin{itemize}
  \item unexplained discontinuities;
  \item overly smooth curves (interpolation?);
  \item cherry-picked time windows;
  \item exaggerated y-axis scales;
  \item missing error bars;
  \item identical noise patterns across experiments.
\end{itemize}

\end{frame}

\begin{frame}[fragile]
\frametitle{Conclusion and Discussion}
\textbf{What to check}:
\begin{itemize}
  \item claims match the results?
  \item limitations acknowledged (or omitted)?
  \item future work realistic or hand-waving?
  \item code/data availability addressed?
\end{itemize}

\textbf{Red flags}:
\begin{itemize}
  \item conclusions broader than evidence;
  \item limitations absent or buried;
  \item critical validations deferred to "future work";
  \item reproducibility not addressed.
\end{itemize}

\end{frame}

\section{Red Flags and Suspicious Practices}

\begin{frame}[fragile]
\frametitle{Suspicious Result Patterns}
\textbf{Warning signs}:

{\tiny
\begin{tabular}{p{3.8cm} p{6.2cm}}
\toprule
\textbf{Pattern} & \textbf{Why suspicious} \\
\midrule
All $p < 0.05$ & Multiple tests produce false positives by chance \\
Improvements only on custom data & Overfitting to domain \\
Monotonic improvement with size & May be cherry-picked curve \\
One metric wins, others don't & Metric gaming \\
Only old baselines & Recent ones might be better \\
\bottomrule
\end{tabular}
}

\end{frame}

\begin{frame}[fragile]
\frametitle{Statistical Red Flags: p-Hacking}
\textbf{p-hacking indicators}:
\begin{itemize}
  \item many tests, only positive ones reported;
  \item $p$-values just below 0.05 (not $< 0.01$);
  \item no multiple comparisons correction;
  \item no error bars or significance tests.
\end{itemize}

\end{frame}

\begin{frame}[fragile]
\frametitle{Statistical Red Flags: Misuse}
\textbf{Underpowered studies}:
\begin{itemize}
  \item small samples without power discussion;
  \item single runs (no replication).
\end{itemize}

\textbf{Misused language}:
\begin{itemize}
  \item "significant" (statistical vs. practical);
  \item "95\% chance" (backwards probability).
\end{itemize}

\end{frame}

\begin{frame}[fragile]
\frametitle{Methodological Red Flags: Leakage}
\textbf{Data leakage}:
\begin{itemize}
  \item hyperparameters tuned on test set;
  \item test set visible during training;
  \item test statistics in preprocessing;
  \item global preprocessing (off-the-shelf).
\end{itemize}

\end{frame}

\begin{frame}[fragile]
\frametitle{Methodological Red Flags: Fairness and Robustness}
\textbf{Unfair comparisons}:
\begin{itemize}
  \item baselines not publicly available or tuned;
  \item different preprocessing for method vs. baselines.
\end{itemize}

\textbf{Suspicious generalization}:
\begin{itemize}
  \item works only on specific domain;
  \item improvements vanish on other datasets;
  \item robustness not tested.
\end{itemize}

\end{frame}

\begin{frame}[fragile]
\frametitle{Visual Red Flags}
Watch for suspicious \textbf{figures and images}:

\textbf{Computer-generated results}:
\begin{itemize}
  \item identical noise patterns in figures claiming independent experiments;
  \item pixel-perfect repetition (copy-paste);
  \item discontinuities that don't match stated methodology;
  \item too smooth or too regular (curve fitting or interpolation on test data).
\end{itemize}

\textbf{Example}: Jan Hendrik Schön's papers had identical noise patterns in plots across supposedly independent experiments.

\textbf{Documentation red flags}:
\begin{itemize}
  \item captions that don't match the figure content;
  \item axis labels missing or inconsistent;
  \item figure quality that suggests poor quality control;
  \item legends that don't match the plotted lines.
\end{itemize}

\end{frame}

\section{Computer Science and Machine Learning Specifics}

\begin{frame}[fragile]
\frametitle{Common ML Pitfalls}
\textbf{Neural networks should report}:
\begin{itemize}
  \item architecture (layers, activations, dropout);
  \item initialization scheme;
  \item optimizer and learning rate;
  \item number of runs (best/mean/median?);
  \item how hyperparameters were chosen.
\end{itemize}

\textbf{Deep learning red flags}:
\begin{itemize}
  \item batch size not reported;
  \item learning rate defaults without justification;
  \item single initialization;
  \item no ablation of architecture components.
\end{itemize}

\end{frame}

\begin{frame}[fragile]
\frametitle{Benchmark Gaming}
\textbf{Dataset overfitting}:
\begin{itemize}
  \item tuning on benchmarks until saturation;
  \item may not generalize to similar data;
  \item compared on same benchmarks where tuned.
\end{itemize}

\textbf{Metric gaming}:
\begin{itemize}
  \item metric doesn't reflect actual goal;
  \item accuracy up, latency down (impractical);
  \item ignores domain constraints (fairness, interpretability).
\end{itemize}

\end{frame}

\begin{frame}[fragile]
\frametitle{Reproducibility Concerns}
\textbf{Red flags}:
\begin{itemize}
  \item code unavailable ("upon request");
  \item hard-coded paths in code;
  \item old or unavailable libraries;
  \item random seeds not reported;
  \item no documentation.
\end{itemize}

\textbf{Green flags}:
\begin{itemize}
  \item GitHub with clear instructions;
  \item environment specified (requirements, Docker);
  \item reproduction scripts;
  \item hyperparameter documentation.
\end{itemize}

\end{frame}

\begin{frame}[fragile]
\frametitle{Transfer Learning and Fine-tuning}
\textbf{Be skeptical}:
\begin{itemize}
  \item ImageNet pre-training on different domains;
  \item fine-tuning on tiny datasets;
  \item no ablation of pre-training benefit;
  \item fine-tuned vs. from-scratch comparisons.
\end{itemize}

\textbf{Standard practice}:
\begin{itemize}
  \item compare pre-trained and from-scratch equally;
  \item tune fine-tuning hyperparameters carefully.
\end{itemize}

\end{frame}

\section{Spotting Questionable Research Practices (QRPs)}

\begin{frame}[fragile]
\frametitle{Common QRPs in Computer Science}
\textbf{Selective reporting}:
\begin{itemize}
  \item report only positive results;
  \item cherry-pick datasets where method wins;
  \item redefine success criteria after seeing data.
\end{itemize}

\textbf{Flexibility in analysis}:
\begin{itemize}
  \item choose aggregation/metric that looks best;
  \item post-hoc grouping of data;
  \item threshold data after seeing results.
\end{itemize}

\textbf{HARKing} (Hypothesizing After Results are Known):
\begin{itemize}
  \item claim results were predicted after seeing them;
  \item frame exploratory as confirmatory.
\end{itemize}

\end{frame}

\begin{frame}[fragile]
\frametitle{Documentation of Questionable Practices}
\textbf{Papers sometimes reveal their own QRPs}:
\begin{itemize}
  \item "we ran this experiment three times and report the best result" (cherry-picking);
  \item "we tried many variants and chose the most stable" (HARKing);
  \item "ablations are in the appendix" (but appendix is missing);
  \item "we optimized hyperparameters for each baseline" vs. "we use default hyperparameters" (inconsistent effort).
\end{itemize}

\end{frame}

\begin{frame}[fragile]
\frametitle{Questions to Ask}
When something feels off:

\begin{enumerate}
  \item Would conclusion hold with reasonable parameter variation?
  \item Were hyperparameters tuned on test set?
  \item Are results consistent across datasets?
  \item Could improvements be implementation differences?
\end{enumerate}

\end{frame}

\section{Constructive Feedback}

\begin{frame}[fragile]
\frametitle{Writing an Effective Review}
\textbf{Structure}:
\begin{enumerate}
  \item Summary (1-2 sentences);
  \item Strengths (2-3 points);
  \item Weaknesses (ordered by importance);
  \item Questions and suggestions;
  \item Recommendation.
\end{enumerate}

\textbf{Tone}:
\begin{itemize}
  \item be specific: point to exact issues;
  \item criticize the work, not authors;
  \item be constructive;
  \item be honest about uncertainty.
\end{itemize}

\end{frame}

\begin{frame}[fragile]
\frametitle{The Reviewer's Checklist}
Before submitting your review:

{\footnotesize
\begin{tabular}{p{6.5cm}p{3.5cm}}
\toprule
\textbf{Question} & \textbf{Checked} \\
\midrule
Is the contribution novel and significant? & \\
Are the claims supported by evidence? & \\
Is the method described clearly enough to reproduce? & \\
Are comparisons fair (same computational budget, tuning effort)? & \\
Are limitations acknowledged? & \\
Are there obvious statistical/methodological errors? & \\
Are there red flags for fabrication or falsification? & \\
Is the paper well-written and organized? & \\
\bottomrule
\end{tabular}
}

\end{frame}

\section{Reviewing for Conferences vs. Journals}

\begin{frame}[fragile]
\frametitle{Conference Reviews}
\textbf{Timeline}: 2-4 weeks.

\textbf{Characteristics}:
\begin{itemize}
  \item \textbf{broad scope}: papers from many areas;
  \item \textbf{quick turnaround}: reviewers are rushed;
  \item \textbf{single-blind or double-blind}: reviewer anonymity varies;
  \item \textbf{lower bar}: acceptance rate 15-30\%, still high volume.
\end{itemize}

\textbf{Your review is brief but must address}:
\begin{itemize}
  \item novelty and significance;
  \item correctness of main claims;
  \item clarity;
  \item recommendation.
\end{itemize}

\end{frame}

\begin{frame}[fragile]
\frametitle{Journal Reviews}
\textbf{Timeline}: 3-6 months (often longer).

\textbf{Characteristics}:
\begin{itemize}
  \item \textbf{specialized scope}: reviewers are experts in the domain;
  \item \textbf{thorough evaluation}: multiple detailed reviews;
  \item \textbf{typically double-blind}: identities hidden;
  \item \textbf{higher bar}: acceptance rate 5-20\%.
\end{itemize}

\textbf{Your review should be comprehensive}:
\begin{itemize}
  \item detailed assessment of methodology;
  \item thorough examination of results;
  \item comparison to related work;
  \item constructive feedback for improvement.
\end{itemize}

\end{frame}

\begin{frame}[fragile]
\frametitle{Open-Science and Preprint Review}
\textbf{Preprints} (arXiv, bioRxiv, medRxiv):
\begin{itemize}
  \item unvetted, rapid dissemination;
  \item \textbf{you can comment publicly} on preprints;
  \item valuable for catching errors before journal acceptance;
  \item community review complements formal peer review.
\end{itemize}

\textbf{Advantages}:
\begin{itemize}
  \item faster feedback to authors;
  \item public accountability (trolling risk, but also prevents editorial bias);
  \item archival record of review process.
\end{itemize}

\end{frame}

\section{Reviewing Your Own Work: Personal Research}

\begin{frame}[fragile]
\frametitle{Reading Papers in Your Field}
As a researcher, you read papers to:
\begin{itemize}
  \item \textbf{keep current}: understand new methods, results, trends;
  \item \textbf{find prior art}: for literature reviews, novelty claims;
  \item \textbf{learn techniques}: understand how to apply methods;
  \item \textbf{evaluate claims}: decide what to build on.
\end{itemize}

Apply the same critical reading to papers you consider building on.

\end{frame}

\begin{frame}[fragile]
\frametitle{Detecting Over-Claimed Results}
\textbf{Common over-claims}:
\begin{itemize}
  \item "state-of-the-art" vs. old baselines;
  \item "efficient" vs. naive implementation;
  \item "robust" on similar distributions;
  \item "solves" one variant of problem.
\end{itemize}

\textbf{Check}:
\begin{itemize}
  \item what exactly was measured?
  \item are baselines standard?
  \item what limitations are acknowledged?
  \item does it work on your data?
\end{itemize}

\end{frame}

\begin{frame}[fragile]
\frametitle{Integrating External Code}
Before using an implementation:

\textbf{Check}:
\begin{itemize}
  \item does code match described algorithm?
  \item results match paper numbers?
  \item well-documented?
  \item reproducibility issues?
\end{itemize}

\end{frame}

\begin{frame}[fragile]
\frametitle{Building a Personal Research Notebook}
As you read papers:
\begin{itemize}
  \item record your assessment of contribution and evidence;
  \item note red flags and questionable claims;
  \item track which papers build on which;
  \item plan validation experiments.
\end{itemize}

This helps you avoid repeating mistakes and find improvement opportunities.

\end{frame}

\section{Common Mistakes Reviewers Make}

\begin{frame}[fragile]
\frametitle{Mistakes to Avoid: Fairness}
\textbf{Being too harsh}:
\begin{itemize}
  \item rejecting papers for incompleteness;
  \item expecting perfection;
  \item criticizing authors instead of work.
\end{itemize}

\textbf{Being too lenient}:
\begin{itemize}
  \item accepting obvious flaws;
  \item not asking for evidence;
  \item following the herd consensus.
\end{itemize}

\end{frame}

\begin{frame}[fragile]
\frametitle{Mistakes to Avoid: Bias and Scope}
\textbf{Over-scoping}:
\begin{itemize}
  \item expecting theory in empirical papers or vice versa;
  \item reviewing outside your expertise.
\end{itemize}

\textbf{Conflicts of interest}:
\begin{itemize}
  \item reviewing papers citing your work;
  \item reviewing competitors;
  \item not declaring conflicts;
  \item promoting your own work.
\end{itemize}

\end{frame}

\begin{frame}[fragile]
\frametitle{Staying Objective}
\textbf{Bias mitigation}:
\begin{itemize}
  \item review the work, not the authors or their affiliations;
  \item be aware of your own biases (field biases, methodological biases);
  \item if you're too close to the topic, consider declining the review;
  \item re-read your review for unfair language before submitting.
\end{itemize}

\textbf{When to decline}:
\begin{itemize}
  \item you lack expertise in the method or domain;
  \item you have a conflict of interest;
  \item you don't have time to do a thorough review;
  \item you dislike the authors or their prior work (this affects objectivity).
\end{itemize}

\end{frame}

\section{Summary and Practical Advice}

\begin{frame}[fragile]
\frametitle{Key Principles}
\begin{enumerate}
  \item \textbf{be a critical reader}: scrutinize methodology, results, and claims;
  \item \textbf{distinguish honest errors from manipulation}: errors are common, malice is rare;
  \item \textbf{check for red flags} in statistics, methodology, and presentation;
  \item \textbf{be constructive}: provide feedback that helps authors improve;
  \item \textbf{stay objective}: review the work, not the authors;
  \item \textbf{know your limits}: don't review outside your expertise.
\end{enumerate}

\end{frame}

\begin{frame}[fragile]
\frametitle{Reviewing Checklist}
{\footnotesize
\begin{enumerate}
  \item Read the abstract and introduction to understand the scope
  \item Identify the main claims and evidence
  \item Check methodology for red flags (data leakage, unfair comparisons, weak baselines)
  \item Assess statistical rigor (error bars, significance tests, multiple comparisons)
  \item Look for cherry-picked results or selective reporting
  \item Verify that conclusion matches the evidence
  \item Check clarity of writing and reproducibility claims
  \item Consider whether the work is novel and significant
  \item Write a fair, constructive review
  \item Declare any conflicts of interest
\end{enumerate}
}

\end{frame}

\begin{frame}[fragile]
\frametitle{Resources}
\textbf{For deeper learning}:
\begin{itemize}
  \item NIH: Guidelines for responsible conduct of research;
  \item Ioannidis (2005): "Why Most Published Research Findings Are False";
  \item Gelman and Loken (2013): "The garden of forking paths";
  \item open science frameworks: Open Science Foundation, FAIR data principles;
  \item domain-specific guidelines: ML reproducibility standards, computational science practice guides.
\end{itemize}

\end{frame}

\begin{frame}[fragile]
\frametitle{Final Thought}
Good peer review is essential to science. As a reviewer, you help maintain research integrity. Be thorough, fair, and constructive — and remember that peer review is ultimately about helping the research community move closer to truth.
\end{frame}




\end{document}
